\documentclass[11pt]{article}

\usepackage{amsmath}
\usepackage{amssymb}
\usepackage{graphicx}
\usepackage{booktabs}
\usepackage{hyperref}
\usepackage{xcolor}
\usepackage{algorithm}
\usepackage{algpseudocode}
\usepackage{multirow}
\usepackage{url}
\usepackage{listings}
\usepackage{geometry}
\usepackage{amsthm}
\usepackage{tikz}
\usepackage{pgfplots}
\pgfplotsset{compat=1.17}
\usepgfplotslibrary{groupplots}

% Evaluation claims in this version are grounded in
% benchmarks/live_audit_recovery.py and the checked-in
% benchmarks/live_audit_recovery_results.json artifact.
\newcommand{\TBD}[1]{\textcolor{red}{\textbf{[TBD: #1]}}}

\newtheorem{theorem}{Theorem}
\newtheorem{lemma}[theorem]{Lemma}
\newtheorem{definition}[theorem]{Definition}

\geometry{margin=1in}

\hypersetup{
  colorlinks=true,
  linkcolor=blue!60!black,
  citecolor=blue!60!black,
  urlcolor=blue!60!black
}

\lstset{
  basicstyle=\ttfamily\small,
  keywordstyle=\color{blue!70!black}\bfseries,
  commentstyle=\color{green!50!black},
  stringstyle=\color{red!60!black},
  breaklines=true,
  frame=single,
  numbers=left,
  numberstyle=\tiny\color{gray},
  language=Python
}

\title{\textbf{TaintFlow: Quantitative Information-Flow Analysis\\for ML Pipeline Leakage Detection}}

\author{%
  Anonymous Authors\\
  \textit{Under Review}
}

\date{}

\begin{document}

\maketitle

%% ============================================================
%% ABSTRACT
%% ============================================================
\begin{abstract}
Data leakage---the inadvertent use of test-set or deployment-time
information during model training---remains one of the most pervasive
yet under-diagnosed pitfalls in machine learning practice.
Existing detection tools rely on syntactic pattern matching or
rule-based heuristics, which limit their coverage to a fixed set of
known anti-patterns and provide no measure of leakage severity.
We present \textsc{TaintFlow}, a research prototype that frames ML
pipeline leakage detection as a \emph{quantitative information-flow}
(QIF) problem.  The repository contains a substantial analysis design
(partition-taint lattice, channel-capacity catalog, and worklist
solver) together with an executable sklearn audit layer that records
live \texttt{fit}, \texttt{fit\_transform}, \texttt{transform}, and
\texttt{score} events via audited wrappers and replay.  In the current
repository snapshot, the strongest grounded evidence comes from a live
benchmark of eight executable sklearn scenarios (four leaky, four
clean) run through this wrapper layer: a simple checker over the
resulting audit logs separates the leaky and clean cases in that small
benchmark, while two representative clean scenarios incur only
near-1x runtime overhead.  A complementary synthetic stress test shows
why structural auditing matters: fitting \texttt{SelectKBest} before
the split on a random-label task inflates test accuracy from 58.3\% to
74.2\%.  We therefore present \textsc{TaintFlow} as a promising,
partially validated audit substrate and QIF-based analysis direction,
rather than as a fully benchmarked end-to-end detector for arbitrary
Python pipelines.
\end{abstract}

%% ============================================================
%% 1  INTRODUCTION
%% ============================================================
\section{Introduction}\label{sec:intro}

Machine learning pipelines are complex software artifacts that
orchestrate data ingestion, feature engineering, model fitting, and
evaluation.  A recurring correctness concern in such pipelines is
\emph{data leakage}: the flow of information from the test (or
hold-out) partition into the training process~\cite{kaufman2012}.
Leakage inflates evaluation metrics and produces models whose
real-world performance falls far short of reported figures.

The problem is widespread.  Kaufman et al.~\cite{kaufman2012}
formalised data leakage in data mining and documented its prevalence
in academic work; practitioners routinely observe that a significant
fraction of public ML notebooks contain at least one leakage
pattern---most commonly fitting a preprocessing step (e.g.,
standardisation, imputation, encoding) on the entire dataset before
performing a train/test split.  Despite its prevalence, leakage is
notoriously difficult to detect manually: the offending data flow may
be buried deep inside library calls, and the resulting metric
inflation can be subtle.

Existing automated approaches fall into two broad categories:
(1)~\emph{syntactic detectors} that match known anti-patterns in
abstract syntax trees; and
(2)~\emph{library-level guards} such as scikit-learn's
\texttt{Pipeline}~\cite{pedregosa2011}, which enforce correct
cross-validation usage within the library's own API.
Each has significant limitations: syntactic detectors miss novel
patterns, and library guards apply only within a single framework.

None of these approaches \emph{quantify} the severity of a detected
leak.  Yet severity matters: a one-bit leak from a binary indicator
may be benign, while a 32-bit leak from a continuous target variable
can render evaluation meaningless.

We propose \textsc{TaintFlow}, a tool that addresses both the
\emph{detection} and \emph{quantification} problems simultaneously by
recasting pipeline leakage as a quantitative information-flow (QIF)
analysis~\cite{smith2009,alvim2020}.  \textsc{TaintFlow} makes the
following contributions:

\begin{enumerate}
  \item A \emph{partition-taint lattice} that tracks how train, test,
    and external data mix through pipeline operations
    (Section~\ref{sec:approach}).
  \item A \emph{channel-capacity bounding} algorithm that computes
    upper bounds on the mutual information leaked, using a catalog of
    80+ scikit-learn and pandas operations
    (Section~\ref{sec:impl}).
  \item A grounded evaluation of the executable sklearn audit layer on
    eight live scenarios (four leaky, four clean), together with a
    synthetic stress test illustrating that structural leakage can
    materially inflate held-out scores even when simpler real-dataset
    examples show no score delta (Section~\ref{sec:eval}).
  \item An open-source implementation with SARIF~\cite{sarif2020}
    output for CI/CD integration (Section~\ref{sec:impl}).
\end{enumerate}

%% ============================================================
%% 2  BACKGROUND
%% ============================================================
\section{Background}\label{sec:background}

\subsection{Data Leakage in ML Pipelines}

A typical supervised-learning pipeline can be modelled as a sequence
of transformations $T_1, T_2, \ldots, T_k$ applied to a dataset
$D = D_{\mathit{tr}} \cup D_{\mathit{te}}$ before a model $M$ is
fitted on $D_{\mathit{tr}}$ and evaluated on $D_{\mathit{te}}$.
Leakage occurs whenever any $T_i$ is \emph{fitted} (i.e., its
parameters are estimated) using samples from $D_{\mathit{te}}$, so
that the transformation applied to $D_{\mathit{tr}}$ implicitly
depends on $D_{\mathit{te}}$.

Common leakage patterns include: (i)~fitting a scaler on the full
dataset before splitting, (ii)~performing feature selection using
labels from all samples, (iii)~target encoding with global
statistics, and (iv)~temporal leakage in time-series data where
future values inform past predictions.

\subsection{Quantitative Information Flow}

Quantitative information flow (QIF)~\cite{smith2009,clark2007}
measures the amount of secret information that an attacker can learn
by observing the public outputs of a system.  Given a secret random
variable $X$ and an observable $Y$, the leakage is typically
characterised by the mutual information $I(X;Y)$ or by the channel
capacity:
\begin{equation}\label{eq:capacity}
  C = \max_{p(x)} I(X;Y).
\end{equation}
An upper bound on $C$ provides a worst-case guarantee: if $C$ is
small, the system leaks little regardless of the input distribution.

\subsection{Abstract Interpretation}

Abstract interpretation~\cite{cousot1977} provides a framework for
computing sound over-approximations of program behaviour.  Given a
concrete semantics operating on a (possibly infinite) domain, one
defines an \emph{abstract domain} connected to the concrete domain
via a Galois connection $(\alpha, \gamma)$.  Computations in the
abstract domain over-approximate the concrete ones, ensuring that any
property proved abstractly holds concretely.

\textsc{TaintFlow} combines QIF and abstract interpretation:
the abstract domain tracks partition membership (a taint), while the
QIF component bounds the capacity of information channels created by
tainted operations.

%% ============================================================
%% 3  APPROACH
%% ============================================================
\section{Approach}\label{sec:approach}

\subsection{Partition-Taint Lattice}

We define a \emph{partition-taint lattice}
$\mathcal{T} = (\mathcal{P}(\{\mathit{tr}, \mathit{te}, \mathit{ext}\})
  \times [0, 64],\; \sqsubseteq)$
where $\mathcal{P}(\{\mathit{tr}, \mathit{te}, \mathit{ext}\})$ is
the powerset of partition labels (training, test, external) and
$[0, 64]$ is a numeric interval representing the maximum number of
bits of information that may have flowed from one partition to
another.

The partial order $\sqsubseteq$ is defined component-wise:
$(S_1, b_1) \sqsubseteq (S_2, b_2)$ iff $S_1 \subseteq S_2$ and
$b_1 \leq b_2$.  The join operator is:
\begin{equation}
  (S_1, b_1) \sqcup (S_2, b_2) = (S_1 \cup S_2,\; \max(b_1, b_2)).
\end{equation}
The bottom element is $(\emptyset, 0)$ and the top is
$(\{\mathit{tr}, \mathit{te}, \mathit{ext}\}, 64)$.

Each variable $v$ in the pipeline is mapped to a taint element
$\tau(v) \in \mathcal{T}$.  After the initial split,
$\tau(D_{\mathit{tr}}) = (\{\mathit{tr}\}, 0)$ and
$\tau(D_{\mathit{te}}) = (\{\mathit{te}\}, 0)$.  A zero-bit count
with a single partition label indicates \emph{no cross-partition
leakage}.

\subsection{PI-DAG Extraction}

\textsc{TaintFlow} constructs a \emph{partition-indexed directed
acyclic graph} (PI-DAG) representing the data-flow dependencies
of a pipeline execution.  Nodes correspond to data objects
(DataFrames, arrays, scalars) and edges to operations.  Each edge is
annotated with:
\begin{itemize}
  \item The \emph{operation type} (e.g., \texttt{fit},
    \texttt{transform}, \texttt{merge}).
  \item The \emph{transfer function} $f_{\mathit{op}}$ that maps
    input taints to output taints.
  \item The \emph{channel capacity bound} $C_{\mathit{op}}$
    contributed by the operation.
\end{itemize}

\subsection{Channel Capacity Bounds}

For each operation in the catalog, \textsc{TaintFlow} provides a
transfer function that computes the output taint and an upper bound
on the bits leaked.  We distinguish three categories:

\paragraph{Partition-preserving operations.}
Operations such as element-wise arithmetic, row-wise filtering, or
column selection that do not mix partition information.  For these,
$C_{\mathit{op}} = 0$ and the partition set is inherited from the
input.

\paragraph{Statistic-aggregating operations.}
Operations such as \texttt{fit} on a \texttt{StandardScaler} that
compute summary statistics (mean, variance).  When the input is a
single partition, $C_{\mathit{op}} = 0$.  When the input mixes
partitions, we bound the capacity by the number of bits in the
computed statistics:
\begin{equation}
  C_{\mathit{fit}} \leq n_{\mathit{params}} \times b_{\mathit{prec}}
\end{equation}
where $n_{\mathit{params}}$ is the number of fitted parameters and
$b_{\mathit{prec}}$ is the floating-point precision (typically 64).

\paragraph{Linear-Gaussian special case.}
For linear transformations of Gaussian data (e.g.,
\texttt{StandardScaler}, \texttt{PCA}), we exploit the closed-form
expression for mutual information:
\begin{equation}\label{eq:gaussian}
  I(X;Y) = \frac{1}{2}\log\det\!\left(I + \Sigma_{X}\Sigma_{N}^{-1}\right)
\end{equation}
where $\Sigma_X$ and $\Sigma_N$ are the signal and noise covariance
matrices, yielding tight bounds.

\subsection{Provenance-Parameterized Channel Capacity}

The key insight of \textsc{TaintFlow} is that the information-theoretic
channel capacity of each operation depends critically on the
\emph{observed provenance}---specifically, the fraction $\rho$ of test
rows in the operation's input.  Consider a \texttt{StandardScaler.fit()}
called on a DataFrame containing $n_{\text{tr}}$ training rows and
$n_{\text{te}}$ test rows.  The fitted \texttt{mean\_} parameter for
feature $j$ is the sample mean:
\begin{equation}
\bar{X}_j = \frac{1}{n}\sum_{i=1}^{n} X_{ij}
\end{equation}
Viewed as a channel from the $n_{\text{te}}$ test rows to this scalar
output, the channel capacity is bounded by:
\begin{equation}\label{eq:cmean}
C_{\text{mean}} \leq \frac{1}{2}\log_2\left(1 + \frac{n_{\text{te}}}{n - n_{\text{te}}}\right) = \frac{1}{2}\log_2\left(1 + \frac{\rho}{1 - \rho}\right) \text{ bits}
\end{equation}
For $\rho = 0.2$ (typical 80/20 split), $C_{\text{mean}} = 0.161$ bits
per feature.  For $\rho = 0.5$, $C_{\text{mean}} = 0.5$ bits.  For
$\rho = 0.01$ (1\% contamination), $C_{\text{mean}} \approx 0.007$ bits,
which is negligible.  This provenance-parameterized approach is what
distinguishes \textsc{TaintFlow} from binary (leakage/no-leakage) tools:
it quantifies severity proportional to the actual contamination level.

For the variance parameter, an analogous bound holds:
\begin{equation}
C_{\text{var}} \leq \frac{1}{2}\log_2\left(1 + \frac{\rho}{1 - \rho}\right) \text{ bits}
\end{equation}

For a \texttt{StandardScaler} with $d$ features, the total channel
capacity of the \texttt{fit} phase is:
\begin{equation}
C_{\text{StandardScaler.fit}} \leq d \cdot (C_{\text{mean}} + C_{\text{var}}) = d \cdot \log_2\left(1 + \frac{\rho}{1 - \rho}\right)
\end{equation}

\subsection{The Data Processing Inequality for Pipeline Composition}

The \emph{data processing inequality} (DPI)~\cite{cover2006} states that
for any Markov chain $X \to Y \to Z$:
\begin{equation}
I(X; Z) \leq \min(I(X; Y), I(Y; Z))
\end{equation}
In an ML pipeline, sequential operations form a Markov chain, so the
DPI guarantees that information can only decrease through lossy stages.
This enables \textsc{TaintFlow} to compose capacity bounds:
for sequential stages $\phi_1, \phi_2, \ldots, \phi_K$, the total
leakage is bounded by:
\begin{equation}
I(D_{\text{te}}; \pi(D)_j) \leq \min_{k=1}^{K} C_{\phi_k}(\rho_k, n_k, d_k)
\end{equation}
The min-cut on the capacity-weighted PI-DAG identifies the
\emph{bottleneck stage}---the operation that is the tightest bound on
information flow---enabling precise attribution.

For parallel branches (e.g., \texttt{ColumnTransformer} processing
different feature subsets independently), the chain rule applies:
\begin{equation}
C_{\text{parallel}} = \sum_{i=1}^{B} C_{\text{branch}_i}
\end{equation}

\subsection{Fit-Transform Decomposition Lemma}

A key insight is that scikit-learn's \texttt{fit}/\texttt{transform}
API naturally decomposes into an information-gathering phase
(\texttt{fit}) and an information-application phase
(\texttt{transform}).

\begin{theorem}[Fit-Transform Decomposition]
Let $\phi$ be a scikit-learn estimator with a \texttt{fit\_transform}
method.  Under the \emph{sufficient-statistic assumption}---that
$\theta = \phi.\text{fit}(D)$ is a sufficient statistic for the
operation $\phi.\text{transform}(x_i, \theta)$---we have:
\begin{equation}
I(D_{\text{te}}; y_j) \leq C_{\text{fit}}(\rho) + b_{\text{input}_j}
\end{equation}
where $C_{\text{fit}}(\rho)$ is the channel capacity of the aggregation
(fit) phase and $b_{\text{input}_j}$ is the existing taint of the
input feature.
\end{theorem}

\begin{proof}[Proof sketch]
Under the sufficient-statistic assumption, the transform output
$y_j = \phi.\text{transform}(x_j, \theta)$ depends on the test
partition $D_{\text{te}}$ only through two paths:
(1)~through $\theta$, since $\theta$ summarises everything
$\text{fit}$ learned from the full dataset, and
(2)~through the input $x_j$ itself, which carries its existing taint
$b_{\text{input}_j}$.
These paths form the Markov chain
$D_{\text{te}} \to (\theta, x_j) \to y_j$.
By the data processing inequality~\cite{cover2006},
$I(D_{\text{te}}; y_j) \leq I(D_{\text{te}}; \theta, x_j)$.
Because $\theta$ and $x_j$ are derived from disjoint operations
(aggregation vs.\ row lookup), we bound the joint channel by the sum:
$I(D_{\text{te}}; \theta, x_j) \leq I(D_{\text{te}}; \theta) +
b_{\text{input}_j} = C_{\text{fit}}(\rho) + b_{\text{input}_j}$.
The inequality is an upper bound and is tight when $\theta$ and
$x_j$ carry independent information about $D_{\text{te}}$.
Note that the bound is conservative when $\text{transform}$ is lossy
(e.g., discretisation), as additional information is lost.
The sufficient-statistic assumption fails for instance-based methods
such as \texttt{KNNImputer}, where the transform output depends on the
full training set rather than a fixed-dimensional summary;
\textsc{TaintFlow} falls back to a conservative full-leak bound for
these operations.
\end{proof}

This decomposition is critical because it allows \textsc{TaintFlow} to
attribute leakage precisely to the \texttt{fit} call (where
cross-partition contamination occurs) rather than the
\texttt{transform} call (which is a deterministic function of the
already-contaminated parameters).

The sufficient-statistic assumption holds for \texttt{StandardScaler}
(mean, variance), \texttt{MinMaxScaler} (min, max), \texttt{PCA}
(covariance matrix), and most linear estimators.  It does \emph{not}
hold for \texttt{KNNImputer} or \texttt{IterativeImputer}, where the
transform output depends on the full training data, not just a
fixed-dimensional statistic.  For these operations, \textsc{TaintFlow}
falls back to the conservative full-leak bound.

\subsection{Worklist Fixpoint}

Taint propagation is computed via a standard worklist fixpoint
algorithm (Algorithm~\ref{alg:worklist}).  The algorithm iterates
until no taint changes, guaranteeing termination because the lattice
has finite height ($2^3 \times 65 = 520$ elements per variable).

\begin{algorithm}[t]
\caption{Worklist fixpoint for taint propagation}\label{alg:worklist}
\begin{algorithmic}[1]
\Require PI-DAG $G = (V, E)$, initial taint map $\tau_0$
\Ensure Final taint map $\tau$
\State $\tau \gets \tau_0$
\State $W \gets$ all nodes with incoming edges
\While{$W \neq \emptyset$}
  \State $v \gets W.\textsc{Pop}()$
  \For{each incoming edge $(u, v, f_{\mathit{op}}, C_{\mathit{op}})$}
    \State $\tau' \gets f_{\mathit{op}}(\tau(u)) \sqcup \tau(v)$
    \If{$\tau' \neq \tau(v)$}
      \State $\tau(v) \gets \tau'$
      \State $W \gets W \cup \textsc{Successors}(v)$
    \EndIf
  \EndFor
\EndWhile
\State \Return $\tau$
\end{algorithmic}
\end{algorithm}

%% ============================================================
%% 4  IMPLEMENTATION
%% ============================================================
\section{Implementation}\label{sec:impl}

\textsc{TaintFlow} is implemented in approximately 54,800 lines of
Python organised into 22 modules under a \texttt{taintflow} package.
We describe its main components below.

\subsection{Architecture}

The codebase is organised as follows:
\begin{itemize}
  \item \texttt{taintflow.core}: The partition-taint lattice (\texttt{lattice.py}),
    type definitions (\texttt{types.py}), configuration (\texttt{config.py}),
    and error hierarchy (\texttt{errors.py}).
  \item \texttt{taintflow.analysis}: The worklist-based fixpoint analyzer
    (\texttt{WorklistAnalyzer}) that propagates taint through the PI-DAG.
  \item \texttt{taintflow.capacity}: Channel capacity models (\texttt{channels.py})
    and the operation catalog (\texttt{catalog.py}) with 80+ transfer functions.
  \item \texttt{taintflow.dag}: PI-DAG data structures (\texttt{nodes.py},
    \texttt{edges.py}), DAG construction (\texttt{builder.py}),
    serialization (\texttt{serialization.py}), and visualization.
  \item \texttt{taintflow.instrument}: Runtime instrumentation via
    \texttt{sys.settrace} (\texttt{tracer.py}), context management
    (\texttt{context.py}), and row provenance tracking (\texttt{provenance.py}).
  \item \texttt{taintflow.integrations}: Framework-specific wrappers for
    scikit-learn (\texttt{sklearn\_interceptor.py}), pandas
    (\texttt{pandas\_interceptor.py}), PyTorch (\texttt{torch\_interceptor.py}),
    and file format support (\texttt{formats.py}).
  \item \texttt{taintflow.attribution}: Min-cut-based leakage attribution
    that identifies bottleneck operations (\texttt{mincut.py}).
  \item \texttt{taintflow.report}: Output formatters for text, JSON, HTML,
    and SARIF formats.
  \item \texttt{taintflow.empirical}: Optional KSG mutual information
    estimation (\texttt{ksg.py}), sliced mutual information estimation
    (\texttt{sliced\_mi.py}), bootstrap confidence intervals
    (\texttt{bootstrap.py}), and permutation tests (\texttt{permutation.py}).
  \item \texttt{taintflow.cli}: Command-line interface with subcommands for
    \texttt{audit}, \texttt{scan}, \texttt{report}, \texttt{compare}, and
    \texttt{config}.
  \item \texttt{taintflow.benchmarks}: Benchmark infrastructure including
    synthetic data generation, scenario definitions, and scoring.
\end{itemize}

\subsection{Runtime Instrumentation}

\textsc{TaintFlow} uses Python's \texttt{sys.settrace} facility to
intercept function calls and returns at runtime.  When a traced
function matches a known operation in the channel catalog,
\textsc{TaintFlow} records the call in the PI-DAG and updates the
taint map.  To reduce overhead, we employ a \emph{selective tracing}
strategy: only functions in whitelisted modules (scikit-learn, pandas,
NumPy) are intercepted; user-defined functions are traced only if
they receive tainted arguments.

The instrumentation layer also provides \emph{audited wrappers} for
common ML objects. For example, \texttt{AuditedStandardScaler} wraps
scikit-learn's \texttt{StandardScaler} and intercepts \texttt{fit()},
\texttt{transform()}, and \texttt{fit\_transform()} to record shapes,
provenance, and fitted parameters in the PI-DAG:

\begin{lstlisting}[caption={Audited estimator usage},label={lst:audited}]
from taintflow.integrations import AuditedStandardScaler
scaler = AuditedStandardScaler()
X_scaled = scaler.fit_transform(X)  # Automatically tracked
# TaintFlow knows: fit saw n rows with test_fraction=rho
\end{lstlisting}

\subsection{PyTorch and File Format Support}

Beyond scikit-learn and pandas, \textsc{TaintFlow} supports:

\paragraph{PyTorch DataLoaders.}  The \texttt{AuditedDataset} and
\texttt{AuditedDataLoader} wrappers track batch-level provenance
through PyTorch's data loading pipeline.  This enables detection of
common deep learning leakage patterns such as computing normalization
statistics (\texttt{transforms.Normalize}) on the full dataset, or
using shared \texttt{Dataset} objects with pre-applied transforms
that were fitted on all data.

\paragraph{File format support.}  Provenance-aware loaders for Parquet
and CSV files (\texttt{load\_parquet}, \texttt{load\_csv}) annotate
DataFrames at ingestion time, enabling TaintFlow to track data lineage
from the point of file loading through the entire pipeline.  A
\texttt{FormatDetector} class auto-detects the file format and applies
appropriate provenance annotations.

\subsection{Channel Catalog}

The channel catalog maps 83 scikit-learn transformers, 14 pandas
operations, and 12 NumPy routines to their transfer functions and
capacity bounds.  Each entry specifies:
\begin{itemize}
  \item Whether the operation is \texttt{fit}-phase,
    \texttt{transform}-phase, or stateless.
  \item The number and type of fitted parameters.
  \item Whether the linear-Gaussian special case applies.
  \item A tightness classification: \textsc{Exact}, \textsc{Near-Tight},
    \textsc{Sound}, or \textsc{Conservative}.
  \item A conservative default bound for unknown operations.
\end{itemize}

Table~\ref{tab:catalog_excerpt} shows representative entries from the
channel catalog with their capacity formulas and tightness factors.

\begin{table}[t]
\centering
\caption{Excerpt from the channel capacity catalog.  $\rho$ is the
  test fraction, $n$ is the sample size, $d$ is dimensionality,
  $c$ is category count.}
\label{tab:catalog_excerpt}
\begin{tabular}{@{}llll@{}}
\toprule
\textbf{Operation} & \textbf{Capacity formula} & \textbf{$\kappa$} & \textbf{Tier} \\
\midrule
\texttt{mean} & $\tfrac{1}{2}\log_2(1+\frac{\rho}{1-\rho})$ & 1 & Tight \\
\texttt{std/var} & $\tfrac{1}{2}\log_2(1+\frac{\rho}{1-\rho})$ & 1 & Tight \\
\texttt{median} & $\log_2(n\rho+1)$ & $O(\log n)$ & Sound \\
PCA ($d$ dims) & $\tfrac{d(d+1)}{2}\cdot C_{\text{cov}}$ & $O(\log d)$ & Sound \\
\texttt{groupby.agg} & $|G|\cdot C_{\text{mean}}(\rho_g)$ & 1 & Tight \\
Target encoding & $H(Y|\text{group})$ & 1 & Sound \\
\texttt{SelectKBest} & $k\cdot C_{\text{mean}}(\rho)$ & $O(1)$ & Sound \\
One-hot encoding & $\log_2(2^c-1)$ & N/A & Conserv. \\
\bottomrule
\end{tabular}
\end{table}

For operations not in the catalog, \textsc{TaintFlow} applies a
conservative \emph{full-leak} assumption: the output taint is the
join of all input taints with $C = 64$ bits per floating-point column.
This ensures soundness while clearly flagging the conservative
treatment in the report.

\subsection{Reporting and Integration}

\textsc{TaintFlow} produces reports in four formats:
(1)~human-readable terminal output with colour-coded severity,
(2)~JSON for programmatic consumption,
(3)~HTML for interactive browsing, and
(4)~SARIF~\cite{sarif2020} for integration with GitHub Code Scanning,
Azure DevOps, and other CI/CD systems.  Each finding includes the source
location, leakage type, estimated bit-count, severity classification,
and a suggested remediation.

The CLI provides six subcommands:
\begin{lstlisting}[language=bash,caption={CLI usage}]
taintflow audit --input pipeline.py        # Full audit
taintflow audit -i pipeline.py -f sarif    # SARIF output
taintflow scan --input pipeline.py         # Quick static scan
taintflow report -i results.json -f html   # Re-render report
taintflow compare --before a.json --after b.json
taintflow config init                      # Generate config
\end{lstlisting}

\subsection{Packaging and Installation}

\textsc{TaintFlow} is organized as a standard Python package with a
\texttt{pyproject.toml} configuration under \texttt{implementation/}.
In the current repository snapshot, the reproducible installation path
is a local editable install:
\begin{lstlisting}[language=bash]
cd implementation
pip install -e .[dev]
\end{lstlisting}

The package requires Python 3.10+ and depends on NumPy, SciPy,
NetworkX, Rich, and Click.  Framework-specific integrations
(scikit-learn, PyTorch, pyarrow) are optional dependencies.

%% ============================================================
%% 5  EVALUATION
%% ============================================================
\section{Evaluation}\label{sec:eval}

We focus the evaluation on executable evidence that is actually present
in this repository snapshot:
(1)~whether the sklearn wrapper layer records enough live structure to
distinguish common pre-split leakage patterns,
(2)~whether structural leakage can matter even when score-based
heuristics are silent, and
(3)~the runtime overhead of the audited wrappers.

\paragraph{Methodology note.}
All numbers in this section come from
\texttt{benchmarks/live\_audit\_recovery.py}, which produces the
checked-in artifact
\texttt{benchmarks/live\_audit\_recovery\_results.json}.  The benchmark
exercises three executable APIs from the current implementation:
\texttt{AuditedEstimator}, \texttt{AuditedPipeline}, and
\texttt{PipelineAuditor}.  Its detection rule is intentionally simple:
flag a scenario whenever an audited \texttt{fit} or
\texttt{fit\_transform} consumes the full dataset before
\texttt{train\_test\_split}.  This validates the wrapper/replay
substrate, not the full \texttt{taintflow audit} CLI.

\subsection{Experiment 1: Live audit replay on executable sklearn pipelines}

\paragraph{Setup.}
We constructed eight executable sklearn scenarios: four leaky and four
clean.  The scenarios cover StandardScaler on iris, PCA on wine,
SelectKBest on breast cancer, and a MinMaxScaler+PCA pipeline on wine.
The leaky variants fit the stateful preprocessing steps on the full
dataset before the split; the clean variants either use
\texttt{AuditedPipeline} directly or replay an ordinary sklearn
\texttt{Pipeline} through \texttt{PipelineAuditor}.  Every scenario is
executed on a real dataset and emits a concrete audit log.

\paragraph{Results.}
Table~\ref{tab:live-recovery} summarizes the benchmark.  The small
benchmark-side checker cleanly separates the four leaky and four clean
cases.  More importantly, the benchmark confirms that the current
repository really can emit stage-level records from live sklearn
execution across both construction-time wrapping
(\texttt{AuditedPipeline}) and replay of pre-existing pipelines
(\texttt{PipelineAuditor}).

\begin{table}[t]
\centering
\caption{Live sklearn audit-replay benchmark generated by
  \texttt{benchmarks/live\_audit\_recovery.py}.}
\label{tab:live-recovery}
\begin{tabular}{@{}lcccc@{}}
\toprule
\textbf{Scenarios} & \textbf{TP} & \textbf{TN} & \textbf{FP} & \textbf{FN} \\
\midrule
8 executable runs & 4 & 4 & 0 & 0 \\
\bottomrule
\end{tabular}
\end{table}

\subsection{Experiment 2: Structural audit versus score deltas}

\paragraph{Setup.}
The same recovery benchmark records held-out scores for four matched
real-dataset pairs (leaky and clean variants of the same pipeline
shape).  We also add a synthetic stress test in which
\texttt{SelectKBest} is fit on a random-label binary classification
task before the split.

\paragraph{Results.}
In the four real-dataset pairs, the leaky and clean variants have
identical held-out scores in this benchmark.  That is a useful result:
it shows why structural leakage auditing cannot be reduced to a simple
``did the score move?'' heuristic.  By contrast, the synthetic
random-label stress test produces a large and reproducible failure
mode: pre-split feature selection increases test accuracy from 58.3\%
to 74.2\%, a 15.8-point absolute gain.  We therefore interpret the
current executable artifact as strongest at \emph{structural auditing},
with the synthetic task serving as a calibration example for why such
auditing matters.

\subsection{Experiment 3: Wrapper overhead}

\paragraph{Setup.}
We measured end-to-end fit-and-score time on two representative clean
pipelines from the recovery benchmark: StandardScaler+LogisticRegression
on iris and PCA+LogisticRegression on wine.  For each scenario we ran
the vanilla sklearn pipeline and the audited variant multiple times and
report the median ratio.

\paragraph{Results.}
Table~\ref{tab:live-overhead} reports the overhead.  The audited runs
remain close to vanilla execution in these small scenarios, suggesting
that the wrapper layer is lightweight enough for interactive debugging
and research use.

\begin{table}[t]
\centering
\caption{Audited-wrapper overhead on representative clean pipelines.}
\label{tab:live-overhead}
\begin{tabular}{@{}lcc@{}}
\toprule
\textbf{Scenario} & \textbf{Vanilla median (s)} & \textbf{Audit ratio} \\
\midrule
\texttt{clean\_scaler\_iris} & 0.00227 & 1.115$\times$ \\
\texttt{clean\_pca\_wine} & 0.01854 & 1.014$\times$ \\
\bottomrule
\end{tabular}
\end{table}

%% ============================================================
%% 6  RELATED WORK
%% ============================================================
\section{Related Work}\label{sec:related}

\paragraph{Leakage in ML pipelines.}
Kaufman et al.~\cite{kaufman2012} formalized data leakage in data
mining and made the case that leakage is both pervasive and subtle.
scikit-learn's \texttt{Pipeline} abstraction~\cite{pedregosa2011}
encourages correct train/test discipline when users structure code in
the intended way, but it does not by itself audit arbitrary scripts or
quantify the severity of a mistake after the fact.  Our current
repository snapshot sits between these positions: it contains a live
wrapper layer that can expose pre-split stateful fits during execution,
and a broader QIF analysis design that aims to turn such structure into
quantitative bounds.

\paragraph{Quantitative information flow.}
Smith~\cite{smith2009} established the foundations of QIF based on
min-entropy, while Alvim et al.~\cite{alvim2020} provide a
comprehensive treatment in their monograph.  Clark, Hunt, and
Malacaria~\cite{clark2007} developed static analyses for quantifying
information flow in imperative programs.  Our work adapts QIF
principles to the specific domain of ML pipelines, where the
``secrets'' are test-partition statistics and the ``channels'' are
pipeline operations.

\paragraph{Abstract interpretation.}
Cousot and Cousot~\cite{cousot1977} introduced the abstract
interpretation framework, which has been widely applied to program
analysis.  Our partition-taint lattice can be viewed as an abstract
domain in this framework, with the fixpoint computation providing
soundness guarantees for the taint propagation.

\paragraph{ML pipeline analysis.}
Several tools analyse ML pipelines for properties other than leakage,
including fairness~\cite{pedregosa2011}, debugging, and
optimisation.  \textsc{TaintFlow}'s PI-DAG extraction could
potentially be extended to support these analyses as well.

%% ============================================================
%% 6.X  MODEL LEAKAGE DETECTION
%% ============================================================
\subsection{Model Leakage Detection}\label{sec:model-leakage}

While \textsc{TaintFlow}'s core contribution targets \emph{data}
leakage---information flowing from test partitions into training---a
complementary threat arises when a trained model itself leaks
information about its training set.  We extend \textsc{TaintFlow}
with an adversarial model-leakage auditor that implements three
membership inference attacks and one model inversion attack.

\paragraph{Membership inference attacks.}
Given black-box access to a target classifier~$f$ trained on
$D_{\mathrm{train}}$, membership inference asks whether a sample~$x$
was in~$D_{\mathrm{train}}$.

\begin{enumerate}
\item \textbf{Shadow model attack}~\cite{shokri2017}: We train
  $k=5$ shadow models on data drawn from the same distribution as
  $D_{\mathrm{train}}$.  The confidence vectors produced by these
  shadow models for their own members and non-members are used to
  train a binary \emph{attack model} (logistic regression), which is
  then evaluated on the target model's members and non-members.

\item \textbf{Threshold attack}~\cite{yeom2018}: A sample~$x$ is
  classified as a member iff its per-sample cross-entropy loss under
  $f$ falls below the mean training loss:
  $\mathbf{1}[\ell(f(x), y) \le \bar{\ell}_{\mathrm{train}}]$.
  This requires only the model's predicted probabilities.

\item \textbf{Label-only attack}~\cite{choquette2021}: When
  confidence scores are unavailable, we perturb each sample with
  Gaussian noise at multiple scales and measure prediction stability.
  Training-set members tend to reside in wider decision regions and
  thus exhibit higher label robustness.
\end{enumerate}

\paragraph{Model inversion.}
We implement gradient-based model inversion~\cite{fredrikson2015},
which reconstructs a class-representative input by maximising the
predicted probability for a target class via numerical gradient
ascent:
$x^* = \arg\max_{x} f_c(x)$, where $f_c$ denotes the predicted
probability for class~$c$.
The quality of the reconstruction is measured by cosine similarity
and a simplified structural similarity index (SSIM) against the
true class centroid.

\paragraph{Privacy budget estimation.}
The connection between membership inference advantage and differential
privacy~\cite{yeom2018} implies $\mathrm{adv} \le e^{\varepsilon}-1$
for $(\varepsilon,\delta)$-DP models.  We therefore report a
\emph{lower bound} on the effective privacy cost:
$\varepsilon \ge \ln(1 + \mathrm{adv})$.

\paragraph{Results.}
Table~\ref{tab:model-leakage} summarises attack success rates on three
standard classifiers trained on a 4\,000-sample synthetic
binary classification task (20 features, 12 informative).

\begin{table}[t]
\centering
\caption{Model leakage audit: membership inference (MI) accuracy and
  advantage, and model inversion cosine similarity.  Higher values
  indicate greater privacy risk.}\label{tab:model-leakage}
\begin{tabular}{@{}llccc@{}}
\toprule
\textbf{Model} & \textbf{Attack} & \textbf{Acc.} & \textbf{Adv.} & \textbf{Cos.\ Sim.} \\
\midrule
\multirow{4}{*}{Logistic Reg.}
  & Shadow model    & 0.483 & 0.033 & N/A \\
  & Threshold       & 0.569 & 0.138 & N/A \\
  & Label-only      & 0.583 & 0.167 & N/A \\
  & Inversion       & N/A   & N/A   & 0.77 \\
\midrule
\multirow{4}{*}{Random Forest}
  & Shadow model    & 0.498 & 0.003 & N/A \\
  & Threshold       & 0.647 & 0.293 & N/A \\
  & Label-only      & 0.357 & 0.286 & N/A \\
  & Inversion       & N/A   & N/A   & $-0.18$ \\
\midrule
\multirow{4}{*}{MLP}
  & Shadow model    & 0.512 & 0.023 & N/A \\
  & Threshold       & 0.631 & 0.262 & N/A \\
  & Label-only      & 0.489 & 0.022 & N/A \\
  & Inversion       & N/A   & N/A   & 0.03 \\
\bottomrule
\end{tabular}
\end{table}

\paragraph{Comparison with TF Privacy and Opacus.}
TensorFlow Privacy's \texttt{membership\_inference\_attack} module
implements threshold and shadow-model attacks but is restricted to
TensorFlow/Keras models.  \textsc{TaintFlow}'s model-leakage auditor
is framework-agnostic---it works with any model exposing
\texttt{predict\_proba}---and additionally implements the label-only
attack of Choquette-Choo et al.~\cite{choquette2021}, which requires
no confidence scores.  Opacus~\cite{opacus2021} provides DP-SGD
training and Rényi DP accounting for PyTorch models but does not
perform empirical leakage measurement.  \textsc{TaintFlow}
complements Opacus by auditing models \emph{after} training---whether
or not differential privacy was applied---and reporting an empirical
lower-bound $\varepsilon$.

\paragraph{When model leakage exceeds data leakage.}
Data leakage inflates evaluation metrics but does not directly expose
individual records.  Model leakage, by contrast, enables an adversary
to infer membership or reconstruct sensitive attributes.  In
privacy-sensitive domains (healthcare, finance), a model with zero
data leakage but high membership inference advantage may pose a
greater deployment risk than one with moderate data leakage.
Integrating both audits provides a more complete privacy posture.

%% ============================================================
%% 7  THREATS TO VALIDITY
%% ============================================================
\section{Threats to Validity}\label{sec:threats}

\emph{Internal validity.}  The live recovery benchmark is intentionally
small: eight sklearn scenarios plus one synthetic stress test.  Its
value is therefore not broad coverage, but rather that every reported
result comes from concrete execution of code in the current repository.

\emph{Construct validity.}  The benchmark-side detector in
Section~\ref{sec:eval} is a simple structural rule over audit logs, not
the full QIF analysis envisioned by the project.  The empirical claims
in this paper should therefore be read as validation of the wrapper and
replay layer, not of the full abstract-interpretation stack.

\emph{External validity.}  Our grounded experiments currently cover
only sklearn-style preprocessing pipelines.  Broader pandas-heavy
feature engineering, temporal pipelines, and the full
\texttt{sys.settrace}-driven path remain future validation targets.

\emph{Score inflation is data dependent.}  In the four paired
real-dataset scenarios, leakage produced no held-out score delta at all;
the synthetic feature-selection task was added precisely because it
makes the structural hazard visible.  This reinforces, rather than
undermines, the need for structural auditing.

%% ============================================================
%% 8  CONCLUSION
%% ============================================================
\section{Conclusion}\label{sec:conclusion}

We presented \textsc{TaintFlow} as a quantitative leakage-auditing
prototype with two distinct layers: a broader QIF-based analysis design
and a presently executable sklearn audit-replay substrate.  In the
current repository snapshot, the strongest grounded evidence comes from
the latter.  The live recovery benchmark shows that audited wrappers
and replay can expose pre-split stateful fits on real sklearn
pipelines, doing so with low overhead in the measured scenarios.

Equally important, the benchmark clarifies \emph{why} this direction is
worth pursuing.  The real-dataset pairs demonstrate that leakage can be
structurally present even when held-out score deltas are absent, while
the synthetic feature-selection stress test shows that the same class
of mistake can dramatically inflate evaluation metrics when conditions
are favorable.

We therefore view the current artifact as a credible foundation for a
stronger future system: one that connects live runtime evidence to the
project's quantitative information-flow machinery in a fully automated
end-to-end CLI.  The immediate next step is to replace the benchmark's
simple structural rule with automatic split inference and analysis over
the extracted execution graph.

%% ============================================================
%% REFERENCES
%% ============================================================
\begin{thebibliography}{10}

\bibitem{kaufman2012}
S.~Kaufman, S.~Rosset, C.~Perlich, and O.~Stitelman,
``Leakage in data mining: Formulation, detection, and avoidance,''
\emph{ACM Trans. Knowl. Discov. Data (TKDD)}, vol.~6, no.~4,
pp.~1--21, 2012.

\bibitem{smith2009}
G.~Smith,
``On the foundations of quantitative information flow,''
in \emph{Proc. 12th Int. Conf. Foundations of Software Science and Computation Structures (FoSSaCS)},
LNCS 5504, pp.~288--302, Springer, 2009.

\bibitem{alvim2020}
M.~S.~Alvim, K.~Chatzikokolakis, A.~McIver, C.~Morgan,
C.~Palamidessi, and G.~Smith,
\emph{The Science of Quantitative Information Flow},
Springer, 2020.

\bibitem{cousot1977}
P.~Cousot and R.~Cousot,
``Abstract interpretation: A unified lattice model for static analysis
of programs by construction or approximation of fixpoints,''
in \emph{Proc. 4th ACM SIGACT-SIGPLAN Symp. Principles of Programming Languages (POPL)},
pp.~238--252, 1977.

\bibitem{clark2007}
D.~Clark, S.~Hunt, and P.~Malacaria,
``A static analysis for quantifying information flow in a simple
imperative language,''
\emph{J. Computer Security}, vol.~15, no.~3, pp.~321--371, 2007.

\bibitem{pedregosa2011}
F.~Pedregosa, G.~Varoquaux, A.~Gramfort, V.~Michel, B.~Thirion,
O.~Grisel, M.~Blondel, P.~Prettenhofer, R.~Weiss, V.~Dubourg,
J.~Vanderplas, A.~Passos, D.~Cournapeau, M.~Brucher, M.~Perrot,
and E.~Duchesnay,
``Scikit-learn: Machine learning in {P}ython,''
\emph{J. Machine Learning Research}, vol.~12, pp.~2825--2830, 2011.

\bibitem{sarif2020}
{OASIS Open},
``Static Analysis Results Interchange Format (SARIF) Version 2.1.0,''
OASIS Standard, 2020.
\url{https://docs.oasis-open.org/sarif/sarif/v2.1.0/sarif-v2.1.0.html}

\bibitem{cover2006}
T.~M.~Cover and J.~A.~Thomas,
\emph{Elements of Information Theory}, 2nd ed.,
Wiley-Interscience, 2006.

\bibitem{openml2013}
J.~Vanschoren, J.~N.~van Rijn, B.~Bischl, and L.~Torgo,
``{OpenML}: Networked science in machine learning,''
\emph{ACM SIGKDD Explorations Newsletter}, vol.~15, no.~2,
pp.~49--60, 2013.

\bibitem{openmlcc18}
B.~Bischl, G.~Casalicchio, M.~Feurer, P.~Gijsbers, F.~Hutter,
M.~Lang, R.~G.~Mantovani, J.~N.~van Rijn, and J.~Vanschoren,
``{OpenML} benchmarking suites,''
in \emph{Proc. NeurIPS 2021 Datasets and Benchmarks Track},
2021.

\bibitem{uci2017}
D.~Dua and C.~Graff,
``{UCI} Machine Learning Repository,''
2017.
\url{https://archive.ics.uci.edu/ml}

\bibitem{shokri2017}
R.~Shokri, M.~Stronati, C.~Song, and V.~Shmatikov,
``Membership inference attacks against machine learning models,''
in \emph{Proc. IEEE Symp. Security and Privacy (S\&P)},
pp.~3--18, 2017.

\bibitem{yeom2018}
S.~Yeom, I.~Giacomelli, M.~Fredrikson, and S.~Jha,
``Privacy risk in machine learning: Analyzing the connection to
overfitting,''
in \emph{Proc. IEEE Computer Security Foundations Symp. (CSF)},
pp.~268--282, 2018.

\bibitem{choquette2021}
C.~A.~Choquette-Choo, F.~Tramer, N.~Carlini, and N.~Papernot,
``Label-only membership inference attacks,''
in \emph{Proc. ICML}, pp.~1964--1974, 2021.

\bibitem{fredrikson2015}
M.~Fredrikson, S.~Jha, and T.~Ristenpart,
``Model inversion attacks that exploit confidence information and
basic countermeasures,''
in \emph{Proc. ACM SIGSAC Conf. Computer and Communications Security (CCS)},
pp.~1322--1333, 2015.

\bibitem{opacus2021}
A.~Yousefpour, I.~Shilov, A.~Sablayrolles, D.~Testuggine,
K.~Prasad, M.~Maber, J.~Lipman, B.~Zong, T.~Mber, L.~Li, and
I.~Mironov,
``{Opacus}: User-friendly differential privacy library in {PyTorch},''
in \emph{NeurIPS Workshop on Privacy-Preserving Machine Learning}, 2021.

\bibitem{goldfeld2021}
Z.~Goldfeld and K.~Greenewald,
``Sliced mutual information: A scalable measure of statistical
dependence,''
in \emph{Proc. Advances in Neural Information Processing Systems (NeurIPS)},
vol.~34, pp.~17567--17578, 2021.

\end{thebibliography}

\end{document}
